\section{Coordinating job recommendations}\label{sec:jobs}
The jobs simulation asks whether coordinating one menu position can recover
placements, and whether that repair improves on an equally coordinated
baseline. It uses Revelio job postings and Census job-seeker profiles, with
modeled preferences and one seat per posting. Each policy shows up to five
feasible postings. The primary design contains 240 markets from August--October
2025 and 160 from the November--December replication, with 24 seekers and
24 postings per market. The same saved \texttt{gpt-6-luna} recommendations,
preferences and feasible pairs are used in the comparisons below.

\begin{table}[ht]\centering\small
\caption{Jobs, primary design. Differences in percentage points. Coverage is
conditional on having a full-list DA partner; placement includes everyone.
The last column compares two coordinated policies and reports point estimates
only: its paired interval is not available in the released aggregates.}
\label{tab:jobs-coordination}
% Generated by jobs/rebuild.py from verified released aggregates.
\begin{tabular}{@{}lrrrr@{}}\toprule
& \multicolumn{2}{c}{GPT minus random} & \multicolumn{2}{c}{Anchored GPT placement minus}\\
Window & Coverage & Placement & Original GPT & Anchored random\\\midrule
Frozen & $+1.6$ & $-15.5$ & $+13.0$ & $-2.81$ \\
Replication & $+0.8$ & $-17.5$ & $+15.9$ & $-2.01$ \\
\bottomrule\end{tabular}

\end{table}

GPT places fewer seekers in both windows. Its coverage advantage is
$1.6$ points in the frozen sample (95\% interval $[0.2,3.1]$) and $0.8$
in replication ($[-1.4,3.3]$). The second interval leaves the coverage
comparison unresolved. As in the rental exercise, this target is a simulated
full-list partner, not observed employment.

The anchor computes a maximum matching on all feasible pairs, preferring
already displayed pairs among equal-cardinality matchings. Each assigned
seeker receives that posting on their menu; a hidden anchor replaces the
last GPT pick, preserving menu length. DA then clears the revised menus.
Proposition~\ref{prop:anchor-capacity} shows that the anchor certifies full
matching capacity when its recipients have nonempty menus, yet a single
edit can still reduce the number placed by DA.
The placement gain is $13.0$ points in the frozen sample and $15.9$ in
replication. Applying the same anchor construction to random menus gives
higher mean placement than anchored GPT: the gaps are $2.81$ and $2.01$
points. Thus the gains relative to the original GPT menus do not establish a
placement advantage over coordinated random exposure.

The released results do not report coverage after anchoring, gross exits,
entrants or individual rank changes. An aggregate gain can therefore coexist
with losses for previous recipients: the menu-edit counts bound possible
gross exits by 459 and 308, respectively, without identifying actual exits
(Proposition~\ref{prop:menu-edit-budget}). Recomputing those endpoints for all four
policies, with paired uncertainty and group outcomes, is the next access
audit; it requires the licensed posting inputs and saved response caches.
No results from that audit are claimed here. Appendix~\ref{app:jobs}
specifies the information sets, existing evidence and pending comparisons.
